% from ega2-1-fr.tex
\oldpage[II]{45}
lorsque $m<0$). On constate comme dans la démonstration de
(\hyperref[II.2.5.13-fr]{2.5.13}) que cette application donne naissance à un
di-homomorphisme de modules
\[
  \eta_f:M'\otimes_{S'}S
  \longrightarrow M'_{(f')}\otimes_{S'_{(f')}}S_{(f)}.
\]
En outre, si, pour $r>0$, on a
$f^r\sum_i(x'_i\otimes s_i)=0$, cela s'écrit aussi
$\sum_i({f'}^rx'_i\otimes s_i)=0$, d'où, par
(\hyperref[0.1.5.4-fr]{0, 1.5.4}),
$\sum_i({f'}^rx'_i/{f'}^{m_i+r})\otimes(s_i/f^{n_i})=0$, c'est-à-dire
$\eta_f(\sum_i x'_i\otimes s_i)=0$, ce qui prouve que $\eta_f$ se factorise
en
\[
  M'\otimes_{S'}S\longrightarrow(M'\otimes_{S'}S)_f
  \xrightarrow{\eta'_f}M'_{(f')}\otimes_{S'_{(f')}}S_{(f)}~;
\]
on vérifie enfin que $\eta'_f$ et $\nu_f$ sont des isomorphismes réciproques
l'un de l'autre.

En particulier, il résulte de (\hyperref[II.2.1.2.1-fr]{2.1.2.1}) que l'on a
un homomorphisme canonique
\[
\label{II.2.8.8.2-fr}
  \Phi^*(\mathcal{O}_{X'}(n))\xrightarrow{\sim}
  \mathcal{O}_X(n)|G(\varphi)
\tag{2.8.8.2}
\]
pour tout $n\in\mathbf{Z}$.
\end{proof}

\begin{env}[2.8.9]
\label{II.2.8.9-fr}
Soient $A$, $A'$ deux anneaux, $\psi:A'\to A$ un homomorphisme d'anneaux,
définissant un morphisme
$\Psi:\operatorname{Spec}(A)\to\operatorname{Spec}(A')$. Soit $S'$ une
$A'$-algèbre graduée à degrés positifs, et posons $S=S'\otimes_{A'}A$, qui
est de façon évidente une $A$-algèbre graduée par les $S'_n\otimes_{A'}A$~;
l'application $\varphi:s'\mapsto s'\otimes1$ est alors un homomorphisme
d'anneaux gradués rendant commutatif le diagramme
(\hyperref[II.2.8.5.1-fr]{2.8.5.1}). Comme ici $S_+$ est le $A$-module
engendré par $\varphi(S'_+)$, on a $G(\varphi)=\operatorname{Proj}(S)=X$~;
d'où, en posant $X'=\operatorname{Proj}(S')$, un diagramme commutatif
\[
\label{II.2.8.9.1-fr}
  \xymatrix{
    X\ar[r]^{\Phi}\ar[d]_p & X'\ar[d]\\
    Y\ar[r]_{\Psi} & Y'
  }
\tag{2.8.9.1}
\]

Soit maintenant $M'$ un $S'$-module gradué, et posons
$M=M'\otimes_{A'}A=M'\otimes_{S'}S$. Dans ces conditions~:
\end{env}

\begin{proposition}[2.8.10]
\label{II.2.8.10-fr}
Le diagramme (\hyperref[II.2.8.9.1-fr]{2.8.9.1}) identifie le schéma $X$ au
produit $X'\times_{Y'}Y$~; en outre, l'homomorphisme canonique
$\nu:\Phi^*(\widetilde{M'})\to\widetilde{M}$
(\hyperref[II.2.8.8-fr]{2.8.8}) est un isomorphisme.
\end{proposition}

\begin{proof}
La première assertion sera démontrée si l'on prouve que pour tout $f'$
homogène dans $S'_+$, en posant $f=\varphi(f')$, les restrictions de $\Phi$
et de $p$ à $D_+(f)$ identifient ce schéma au produit
$D_+(f')\times_{Y'}Y$ (\hyperref[I.3.2.6.2-fr]{I, 3.2.6.2})~; autrement
dit, il suffit (\hyperref[I.3.2.2-fr]{I, 3.2.2}) de prouver que $S_{(f)}$
s'identifie canoniquement à $S'_{(f')}\otimes_{A'}A$, ce qui est immédiat
en vertu de l'existence de l'isomorphisme canonique
$S_f\xrightarrow{\sim}S'_{f'}\otimes_{A'}A$ conservant les degrés
(\hyperref[0.1.5.4-fr]{0, 1.5.4}). La seconde assertion résulte de ce que
$M'_{(f')}\otimes_{S'_{(f')}}S_{(f)}$ s'identifie d'après ce qui précède à
$M'_{(f')}\otimes_{A'}A$, et ce dernier à $M_{(f)}$, puisque $M_f$
s'identifie canoniquement à $M'_{f'}\otimes_{A'}A$ par un isomorphisme
conservant les degrés.
\end{proof}

\begin{corollary}[2.8.11]
\label{II.2.8.11-fr}
Pour tout entier $n\in\mathbf{Z}$, $\widetilde{M(n)}$ s'identifie à
$\Phi^*(\widetilde{M'(n)})=\widetilde{M'(n)}\otimes_{Y'}\mathcal{O}_Y$~;
en particulier
$\mathcal{O}_X(n)=\Phi^*(\mathcal{O}_{X'}(n))
=\mathcal{O}_{X'}(n)\otimes_{Y'}\mathcal{O}_Y$.
\end{corollary}

\begin{proof}
Cela résulte de (\hyperref[II.2.8.10-fr]{2.8.10}) et de
(\hyperref[II.2.5.15-fr]{2.5.15}).
\end{proof}

